\begin{thebibliography}{99}
\bibitem{chernoff} H. Chernoff, ``Sequential design of experiments,'' \emph{Ann. Math. Statist.}, vol. 30, no. 3, pp. 755--770, 1959.
\bibitem{naghshvar} M. Naghshvar and T. Javidi, ``Active sequential hypothesis testing,'' \emph{Ann. Statist.}, vol. 41, no. 6, pp. 2703--2738, 2013, doi: 10.1214/13-AOS1144.
\bibitem{atkinson} A. C. Atkinson and V. V. Fedorov, ``The design of experiments for discriminating between two rival models,'' \emph{Biometrika}, vol. 62, no. 1, pp. 57--70, 1975, doi: 10.1093/biomet/62.1.57.
\bibitem{burnashev} M. V. Burnashev, ``On the minimax detection of an inaccurately known signal in a white Gaussian noise background,'' \emph{Theory Probab. Appl.}, vol. 24, no. 1, pp. 107--119, 1979, doi: 10.1137/1124008.
\bibitem{gjn} A. Goldenshluger, A. Juditsky, and A. Nemirovski, ``Hypothesis testing by convex optimization,'' \emph{Electron. J. Statist.}, vol. 9, no. 2, pp. 1645--1712, 2015, doi: 10.1214/15-EJS1054.
\bibitem{basseville} M. Basseville and I. V. Nikiforov, \emph{Detection of Abrupt Changes: Theory and Application}. Englewood Cliffs, NJ, USA: Prentice Hall, 1993.
\bibitem{zhang2002} X. Zhang, M. M. Polycarpou, and T. Parisini, ``A robust detection and isolation scheme for abrupt and incipient faults in nonlinear systems,'' \emph{IEEE Trans. Autom. Control}, vol. 47, no. 4, pp. 576--593, 2002, doi: 10.1109/9.995036.
\bibitem{nik98} R. Nikoukhah, ``Guaranteed active failure detection and isolation for linear dynamical systems,'' \emph{Automatica}, vol. 34, no. 11, pp. 1345--1358, 1998.
\bibitem{campbell} S. L. Campbell and R. Nikoukhah, \emph{Auxiliary Signal Design for Failure Detection}. Princeton, NJ, USA: Princeton University Press, 2004.
\bibitem{scott} J. K. Scott, R. Findeisen, R. D. Braatz, and D. M. Raimondo, ``Input design for guaranteed fault diagnosis using zonotopes,'' \emph{Automatica}, vol. 50, no. 6, pp. 1580--1589, 2014, doi: 10.1016/j.automatica.2014.03.016.
\bibitem{blanchini} F. Blanchini, D. Casagrande, G. Giordano, S. Miani, S. Olaru, and V. Reppa, ``Active fault isolation: A duality-based approach via convex programming,'' \emph{SIAM J. Control Optim.}, vol. 55, no. 3, pp. 1619--1640, 2017, doi: 10.1137/15M1046046.
\bibitem{heirung} T. A. N. Heirung and A. Mesbah, ``Input design for active fault diagnosis,'' \emph{Annu. Rev. Control}, vol. 47, pp. 35--50, 2019, doi: 10.1016/j.arcontrol.2019.03.002.
\bibitem{simandl} M. \v{S}imandl and I. Pun\v{c}och\'{a}\v{r}, ``Active fault detection and control: Unified formulation and optimal design,'' \emph{Automatica}, vol. 45, pp. 2052--2059, 2009.
\bibitem{ashari} A. Esna Ashari, R. Nikoukhah, and S. L. Campbell, ``Effects of feedback on active fault detection,'' \emph{Automatica}, vol. 48, pp. 866--872, 2012.
\bibitem{coster} D. C. Coster and C.-S. Cheng, ``Minimum cost trend-free run orders of fractional factorial designs,'' \emph{Ann. Statist.}, vol. 16, no. 3, pp. 1188--1205, 1988.
\bibitem{cheng} C.-S. Cheng and D. M. Steinberg, ``Trend robust two-level factorial designs,'' \emph{Biometrika}, vol. 78, no. 2, pp. 325--336, 1991.
\bibitem{bickel1} P. J. Bickel and A. M. Herzberg, ``Robustness of design against autocorrelation in time I: Asymptotic theory, optimality for location and linear regression,'' \emph{Ann. Statist.}, vol. 7, no. 1, pp. 77--95, 1979.
\bibitem{bickel2} P. J. Bickel, A. M. Herzberg, and M. F. Schilling, ``Robustness of design against autocorrelation in time II: Optimality, theoretical and numerical results for the first-order autoregressive process,'' \emph{J. Amer. Statist. Assoc.}, vol. 76, no. 376, pp. 870--877, 1981.
\bibitem{schieder} R. Schieder and C. Kramer, ``Optimization of heterodyne observations using Allan variance measurements,'' \emph{Astron. Astrophys.}, vol. 373, pp. 746--756, 2001.
\bibitem{ossenkopf} V. Ossenkopf, ``The stability of spectroscopic instruments: A unified Allan variance computation scheme,'' \emph{Astron. Astrophys.}, vol. 479, pp. 915--926, 2008.
\bibitem{vitus} M. P. Vitus, W. Zhang, A. Abate, J. Hu, and C. J. Tomlin, ``On efficient sensor scheduling for linear dynamical systems,'' \emph{Automatica}, vol. 48, no. 10, pp. 2482--2493, 2012.
\bibitem{zhao} L. Zhao, W. Zhang, J. Hu, A. Abate, and C. J. Tomlin, ``On the optimal solutions of the infinite-horizon linear sensor scheduling problem,'' \emph{IEEE Trans. Autom. Control}, vol. 59, no. 10, pp. 2825--2830, 2014.
\bibitem{banerjee} T. Banerjee and V. V. Veeravalli, ``Data-efficient quickest change detection with on--off observation control,'' \emph{Sequential Analysis}, vol. 31, no. 1, pp. 40--77, 2012.
\bibitem{lautay} T. S. Lau and W. P. Tay, ``Asymptotically optimal sampling policy for quickest change detection with observation-switching cost,'' \emph{IEEE Trans. Signal Process.}, vol. 69, pp. 1332--1346, 2021, doi: 10.1109/TSP.2021.3057258.
\bibitem{lubenia} P. V. N. Lubenia and T. Banerjee, ``Quickest change detection with cost-constrained experiment design,'' arXiv:2509.14186v2, 2025, preprint.
\bibitem{liang} Y. Liang, A. G. Tartakovsky, and V. V. Veeravalli, ``Quickest change detection with non-stationary post-change observations,'' \emph{IEEE Trans. Inf. Theory}, vol. 69, no. 5, pp. 3400--3414, 2023, doi: 10.1109/TIT.2022.3230583.
\bibitem{hou} Y. Hou, Y. Oleyaeimotlagh, R. Mishra, H. Bidkhori, and T. Banerjee, ``Robust quickest change detection in nonstationary processes,'' \emph{Sequential Analysis}, vol. 43, no. 3, pp. 275--300, 2024, doi: 10.1080/07474946.2024.2356555.
\bibitem{deluca} A. De Luca and R. Mattone, ``Actuator failure detection and isolation using generalized momenta,'' in \emph{Proc. IEEE Int. Conf. Robot. Autom.}, vol. 1, 2003, pp. 634--639, doi: 10.1109/ROBOT.2003.1241665.
\bibitem{haddadin} S. Haddadin, A. De Luca, and A. Albu-Sch\"affer, ``Robot collisions: A survey on detection, isolation, and identification,'' \emph{IEEE Trans. Robot.}, vol. 33, no. 6, pp. 1292--1312, 2017, doi: 10.1109/TRO.2017.2723903.
\bibitem{lynch} K. M. Lynch and F. C. Park, \emph{Modern Robotics: Mechanics, Planning, and Control}. Cambridge, U.K.: Cambridge University Press, 2017, chs. 8 and 11.
\end{thebibliography}
